\section{Compute available for free}
\label{sec:compute}

The hardware decides which models are candidates at all, so it is treated as a design variable, not a
given. Free-tier limits change often; the numbers below were collected in September 2026 from
provider pages and 2026 guides~\cite{kagglegpu,colabfaq,lightningfree} and must be re-checked on
sign-up.

\begin{table}[H]
\centering\small
\begin{tabularx}{\textwidth}{@{}L{2.7cm}L{3.4cm}L{3.3cm}Y@{}}
\toprule
\textbf{Option} & \textbf{Hardware} & \textbf{Limits} & \textbf{Best use in this project} \\
\midrule
Laptop (ours) & i5-6300U, 2C/4T, AVX2 (no AVX-512 / VNNI), 15\,GB RAM, $\sim$2–5\,GB free disk & None; always available & The ``edge'' configuration. Every measurement in Section~\ref{sec:measure}. Final demo client (microphone and loudspeaker). \\
Kaggle Notebooks & 2$\times$T4 (16\,GB each) or 1$\times$P100 & $\sim$30 GPU-h/week; sessions capped (9–12\,h); internet after phone verification & Main GPU benchmark machine: the same notebook reruns every experiment on the same hardware. \\
Google Colab (free) & T4 16\,GB (when available) & Dynamic quota; idle timeout; $\le$12\,h sessions & Second GPU, OmniVoice demo notebook, quick checks. \\
Lightning AI Studio (free tier) & T4/L4 via monthly credits; 1 persistent CPU studio & Monthly credit allowance (sources report 22–80 GPU-h/month) & Persistent environment. Exposing a port for a live demo is simpler than on Kaggle or Colab. \\
Hugging Face Spaces & CPU basic (free); ZeroGPU needs PRO to host & Cold starts; queueing & Public demo of the CPU pipeline only. \\
Oracle Cloud Always Free & 4 Arm Ampere OCPU, 24\,GB RAM & Always on & Optional always-on CPU server (llama.cpp performs well on Arm NEON). \\
AWS SageMaker Studio Lab & – & Reported closed to new users in July 2026 & Not relied on. \\
\bottomrule
\end{tabularx}
\caption{Free compute options. A T4 (Turing, 65 TFLOPS fp16, 320\,GB/s, no bf16) is the realistic
GPU target: all fp16/int8 paths must work without bf16.}
\label{tab:compute}
\end{table}

\paragraph{The network is part of \ttfa.} If the models run in a notebook in a US data centre and the
microphone is in Addis Ababa, every turn pays at least one round trip for the user's audio and one
for the response audio, plus any tunnel overhead (cloudflared, ngrok). This must be measured
(experiment E-SYS-4). The design rule is to \textbf{co-locate all models on one machine} and to stream
audio \emph{while the user speaks}, so that uploading the user's audio overlaps with speaking and
only the response direction costs latency.

\paragraph{Two target configurations.} We therefore optimise two configurations and report both:
\begin{enumerate}
  \item \textbf{CPU-only (edge)}: the laptop, no network. This is the hardest case and the most honest one,
  and everything in Section~\ref{sec:measure} was measured here.
  \item \textbf{Single T4 (server)}: all models on one Kaggle or Colab T4, with the laptop as a thin
  audio client. This is the configuration most likely to reach the ``Good'' tier.
\end{enumerate}
